\documentclass[11pt]{article}
\usepackage[a4paper,margin=23mm]{geometry}
\usepackage{amsmath,amssymb,amsthm,mathtools,bm}
\usepackage[T1]{fontenc}
\usepackage{lmodern,microtype,booktabs}
\usepackage[colorlinks=true,linkcolor=blue,citecolor=blue,urlcolor=blue]{hyperref}
\newcommand{\EE}{\mathbb E}
\newcommand{\NN}{\mathcal N}
\newcommand{\RR}{\mathcal R}
\newcommand{\DD}{\mathfrak D}
\newcommand{\WW}{\mathcal W}
\newcommand{\tr}{\operatorname{tr}}
\newcommand{\PP}{\mathcal P_h}
\newcommand{\thN}{\theta(\NN)}
\newcommand{\rhat}{\widehat r}
\newcommand{\ehat}{\widehat\eta}
\newtheorem{proposition}{Proposition}
\title{A rank-marked series for the neutral memory field\\
\large Actual propagators, Jacobi contraction errors, and finite-time control}
\author{Alan Day \and Chengheng\thanks{AI research assistant (OpenAI).}}
\date{7 November 2026\quad Local research draft}
\begin{document}
\maketitle
\begin{abstract}
We give a uniformly summable representation of the neutral component of
the common-coupling complex DSSYK memory field on every fixed real-time
interval. Its inputs are highest Gaussian-chaos fields with one output-rank
deficit mark. A forced recurrence retains all lower-chaos contractions.
The main additional estimate is uniform integrability of the rank-deficit
square for the actual exponential, obtained from fixed-degree sixth moments
and a graph-norm Taylor tail. This identifies a finite-time computational
interface; the limiting marked joint spectral measure remains to be computed.
\end{abstract}

\section{Model and the result to be represented}
Use standard CAR, overlapping complex Gaussian $p$-body couplings, and the
same original couplings in every charge sector. The state is the normalized
annealed full trace at $\beta=\mu=0$, with weights
$b_m=2^{-N}\binom Nm$. We take even $N$, $p\to\infty$,
$\lambda_N=p^2/N\to\lambda>0$, fixed $J_s$, $\hbar=1$, and string real time.
In sector $m$ write $h=H_m=\sum_\gamma g_\gamma B_\gamma$,
$\sum_\gamma B_\gamma^2=C_m I$; $C_m\le M$ uniformly.
Norms mean $\|X\|_{2,m}^2=\EE\tr_m X^\dagger X$.
An independent copy $u$ defines $\Phi(X)=\EE_u[uXu]$.
The exact energy pinching is $\PP X=\sum_e P_eXP_e$, including
whole degenerate blocks.

Put $\alpha_m=(1-p/m)_+$ and $\epsilon_m=1-\alpha_m^2$ (zero sectors
are interpreted directly). The preceding two-component reduction gives
$\mathsf F=A+B+o_{L^2}(1)$, where the neutral field is
\begin{equation}
 B_m(s)=(1-\PP)m\epsilon_m e^{-ihs/2}
                 \Phi(e^{i\alpha_mhs})e^{-ihs/2}.\label{eq:B}
\end{equation}
Here $A$ is the charged transfer defect from the previous stage
\cite{Two,Uniform}. We retain $\alpha_m$ and $m\epsilon_m/p$ exactly.
Our new representation of $B$ is stated on page~\pageref{result};
its construction uses finite-$N$ coefficients throughout.

\newpage
\section{Exact Jacobi rearrangement, before taking a limit}
In a broad central sector $N/4\le m\le3N/4$, suppress $m$ and set
\begin{align}
 \mu(r)&=\prod_{\ell=0}^{p-1}
 \left(1-\frac r{m-\ell}\right)_+
 \left(1-\frac r{N-m+p-\ell}\right)_+,\qquad \theta_j=\mu(jp),\\
 \DD&=p\NN-\RR,\qquad Z_j=\WW_j(h^j),\qquad U_j=\DD Z_j.
\end{align}
$\NN$ counts coupling Gaussian chaos; $\RR$ is the output $U(N)$
rank. On equivariant functions $\DD\ge0$, and $\Phi/C=\mu(\RR)$.
These operators commute with one another, but not generally with
samplewise multiplication by $h$. The positive-part extension defines
$\mu(jp)$ even beyond the physical rank range.

Let $\lambda_N\in[\lambda_-,\lambda_+]$ and $p\le N/16$. The previous
Jacobi construction \cite{Uniform} supplies
\begin{align}
 a_j&=\sum_{v<j}\theta_v,\quad 1\le a_j\le A\ (j\ge1),
 \quad A=(1-e^{-3\lambda_-})^{-1},\\
 Q_0&=1,\quad Q_1=x,\quad Q_{j+1}=xQ_j-Ca_jQ_{j-1},\\
 E_j&=Z_j-Q_j(h),\qquad \eta_j=pE_j.
\end{align}
Here $a_0=0$ and the three-term recurrence starts at $j=1$.
$E_j$ has Gaussian degree at most $j-2$. The scalar coefficients of
the auxiliary bounded Jacobi operator satisfy the entire identity
and the bounds
\begin{align}
 e^{is h}&=\sum_{j\ge0}c_j(s)Q_j(h),\qquad
 |c_j(s)|\le e^{2\sqrt{AM}S}\frac{(2\sqrt A S)^j}{j!},\label{eq:c}\\
 \|E_j\|_2&\le\frac pN L K_*^j\sqrt{j!}\,C^{j/2},\qquad
 K_*=4+A,\label{eq:E}\\
 L&=12\sqrt{2\cdot10^6}(\sqrt{\lambda_+}+\lambda_+),\qquad |s|\le S.
\end{align}
This is an expansion on the actual Gaussian spectrum, not just on
the compact support of an auxiliary measure.

Since $\thN Z_j=\theta_jZ_j$, exact algebra gives
\begin{equation}
 \thN Q_j(h)=\theta_jQ_j(h)+(\theta_j-\thN)E_j.
\end{equation}
The sum $\sum c_j\theta_jQ_j(h)$ is a samplewise energy function.
All series converge absolutely in annealed $L^2$. Therefore
\begin{equation}
 \boxed{\quad
 \frac pC(1-\PP)\Phi(e^{ish})=(1-\PP)
 \left\{p[\mu(\RR)-\thN]e^{ish}
 +\sum_{j\ge0}c_j(s)(\theta_j-\thN)\eta_j\right\}.
 \quad}\label{eq:exact}
\end{equation}
The second sum cannot be omitted: $pE_j$ is order one on this scale.

\newpage
\section{Fixed-degree higher moments from the existing rank law}
The positive rank-product law and the sixth-moment certificate of
\cite{Central,Rank} extend to any fixed number of rank steps.
For independent pure ranks $a,r$, let the output rank be $t$ and
$S=a+r-t$. Its Racah identity determines moments through
\begin{align}
 \EE\chi_t(n)&=\chi_a(n)\chi_r(n),\qquad x_t=t(N+1-t),\\
 F_n&=\frac{\chi_a(n)\chi_r(n)-\sum_{j<n}c_{nj}F_j}{c_{nn}},
 \quad F_j=\EE\prod_{\ell<j}(x_t-x_\ell).
\end{align}
Here $c_{nj}$ are the coefficients of $\chi_t(n)$ in that monic
polynomial basis; their exact expression and the symbolic certificate
are in \cite{Central}.

The certificate cancels 21 negative Laurent coefficients through
order six over the full polynomial ring $\mathbb Q[A_0,B_0,z]$, with
$A_0=a/\sqrt N$, $B_0=r/\sqrt N$, $z=(m-N/2)/\sqrt N$.
Its denominators have nonzero limiting constants on compact sets.
For any fixed $K,X$, it consequently applies to
\begin{equation}
 a\le p,\quad r\le Kp,\quad |m-N/2|\le Xp.
\end{equation}
Indeed $x_{a+r}-x_t=S[N+1-2(a+r)+S]$, whose bracket is at least
$N/2$ for sufficiently large $N$ at fixed $K$. Hence $\EE S^6$ is
uniformly bounded. This extension changes the compact set of the
old certificate, not the physical interaction size $p$.

For a Hamiltonian rank step define $\delta=p+r-t=(p-a)+S$.
The input rank deficit $p-a$ has the previously proved factorial tail.
Mixture over $a$ and $(x+y)^6\le32(x^6+y^6)$ give finite constants
\begin{equation}
 M_{K,6}=\sup_{r\le Kp}\EE_r\delta^6<\infty,\qquad
 S_{K,4}=\sup_{r\le Kp}\EE_r\delta^4\le M_{K,6}^{2/3}.\label{eq:step6}
\end{equation}
These constants may depend on $K,X,\lambda_\pm$.
The positive Markov chain of $k$ independent inputs has
\begin{equation}
 D_k=kp-r_k=\sum_{j<k}\delta_j,\qquad
 \EE D_k^6\le k^5\sum_{j<k}M_{j,6}.\label{eq:fixed6}
\end{equation}
Gaussian symmetrization is a contraction and commutes with output
rank, so
\begin{equation}
 \|\DD^aZ_k\|_2^2\le k!C^k\EE D_k^{2a},\qquad a=1,2,3.
\end{equation}
This fixed-$k$ conclusion alone does not control the actual exponential.
The next page supplies the degree tail before either cutoff is removed.

\newpage
\section{A graph-norm tail for the actual exponential}
Let $\mathcal T$ denote left multiplication by $h$. On degree at most $r$,
the all-degree first-commutator bound of \cite{Rank} reads
\begin{align}
 \|\mathcal T\psi\|_2&\le L_H\sqrt{r+1}\|\psi\|_2,
 \quad L_H=2\sqrt M,\\
 \|[\DD,\mathcal T]\psi\|_2&\le K_0(r+1)^{3/2}\|\psi\|_2,
 \quad K_0=2000\sqrt{M(\lambda_++\lambda_+^2)}.
\end{align}
These hold throughout the broad central interval.
Commuting $\DD$ through $h^n$ inductively gives
\begin{equation}
 b_n:=\|\DD h^n\|_2\le
 \frac{K_0n(n+1)}2 L_H^{n-1}\sqrt{n!}.
\end{equation}
For the ordinary Taylor polynomial $P_K(s,h)$, closedness of $\DD$
and the convergent Gaussian moment series imply
\begin{align}
 \sup_{|s|\le S}\|\DD(e^{ish}-P_K(s,h))\|_2&\le T_K(S),\\
 T_K(S)&=\frac{K_0}2\sum_{n>K}
 \frac{n(n+1)S^nL_H^{n-1}}{\sqrt{n!}}\longrightarrow0.\label{eq:graph}
\end{align}
In a compact central charge window, rank orthogonality, creation
symmetrization and \eqref{eq:step6} also give
\begin{equation}
 \|[\DD,[\DD,\mathcal T]]\psi_{\le K}\|_2
 \le J_K\|\psi\|_2,\qquad J_K=2\sqrt{M(K+1)S_{K,4}}.
\end{equation}
The annihilation block is the adjoint of the creation block; no
reverse transition probability is assumed. Define $d_0=0$ and
\begin{align}
 d_n={}&L_H\sqrt n\,d_{n-1}+2K_0n^{3/2}b_{n-1}
       +J_{n-1}L_H^{n-1}\sqrt{(n-1)!},\\
 C_K(S)&=\sum_{n\le K}S^nd_n/n! .
\end{align}
The identity $\DD^2\mathcal T=\mathcal T\DD^2+
2[\DD,\mathcal T]\DD+[\DD,[\DD,\mathcal T]]$ proves
$\|\DD^2h^n\|_2\le d_n$. Thus
\begin{equation}
 \boxed{\quad\sup_{|s|\le S}
 \|\DD1_{\DD>R}e^{ish}\|_2\le T_K(S)+C_K(S)/R.\quad}\label{eq:UI}
\end{equation}
Choose $K$ first, then $R$. This is uniform integrability of the
$\DD^2$ spectral weight of the complete exponential.

For fixed $k,d$, differentiating the finite product for $\mu$ gives
$p[\mu(kp-d)-\mu(kp)]-4\lambda_N\theta_kd\to0$.
Globally its first term is at most $12\lambda_Nd$, so the absolute
combined multiplier is bounded by $16\lambda_+\DD$.
Degree cutoff \eqref{eq:graph} and rank cutoff \eqref{eq:UI} therefore prove
\begin{equation}
 p[\mu(\RR)-\thN]e^{ish}-4\lambda_N\thN\DD e^{ish}
 \longrightarrow0\quad\hbox{in }L^2,\label{eq:slope}
\end{equation}
uniformly on compact real times and compact central charge windows.

\newpage
\section{Generating every contraction correction from one marked family}
Use $\diamond$ for the ordinary coupling-Gaussian Wick product with
ordered matrix coefficients, not the spectral $q$-Wick product.
Differentiating the homogeneous polynomial before projecting gives
the exact annihilation remainder
\begin{align}
 r_j&=\sum_{v<j}[(\Phi-C\theta_v)Z_v]\diamond Z_{j-1-v},\label{eq:r}\\
 \eta_{j+1}&=h\eta_j-Ca_j\eta_{j-1}-pr_j,
 \qquad\eta_0=\eta_1=\eta_2=0.\label{eq:eta}
\end{align}
For a homogeneous degree-$v$ tensor $V_v$, right creation by $h$
repeated $w$ times proves
\begin{equation}
 \|V_v\diamond Z_w\|_2\le
 \sqrt{(v+w)!/v!}\,C^{w/2}\|V_v\|_2.\label{eq:wickbound}
\end{equation}
One sums $\sum B_\gamma^2=CI$ before symmetrization at each step.
This estimate does not assert a bound for arbitrary products of two
matrix-valued Gaussian functions.

Define the finite-parameter source and its forced recurrence by
\begin{align}
 \rhat_j&=4\lambda_NC\sum_{v<j}\theta_v
                          (U_v\diamond Z_{j-1-v}),\label{eq:rhat}\\
 \ehat_0&=\ehat_1=0,\qquad
 \ehat_{j+1}=h\ehat_j-Ca_j\ehat_{j-1}-\rhat_j.\label{eq:etahat}
\end{align}
$U_0=0$ also gives $\ehat_2=0$. At fixed $j$, the slope calculation
on $Z_v$, controlled by \eqref{eq:fixed6}, yields
$\|pr_j-\rhat_j\|_2\to0$.
Every error has bounded Gaussian degree, so multiplication by $h$
is bounded on it. Induction in \eqref{eq:eta} gives
\begin{equation}
 \|\eta_j-\ehat_j\|_2\longrightarrow0\qquad(j\text{ fixed}).\label{eq:etalim}
\end{equation}

The all-degree second-moment bound already established in \cite{Rank} is
\begin{equation}
 \|U_v\|_2\le\sqrt{v!}\,C^{v/2}\sqrt{K_v},\qquad
 K_v=2\cdot10^6(\lambda_Nv^2+\lambda_N^2v^4).
\end{equation}
Set $q_*=e^{-3\lambda_-}$. Since $\theta_v\le q_*^v$, \eqref{eq:wickbound}
implies $\|\rhat_j\|_2\le\widehat L C^{(j+1)/2}\sqrt{(j-1)!}$, where
\begin{equation}
 \widehat L=4\lambda_+\sqrt{2\cdot10^6}
 \left[\frac{\sqrt{\lambda_+}q_*}{(1-q_*)^2}
 +\frac{\lambda_+q_*(1+q_*)}{(1-q_*)^3}\right].
\end{equation}
Recurrence \eqref{eq:etahat} now gives the uniform estimate
\begin{equation}
 \|\ehat_j\|_2\le\widehat L K_*^j C^{j/2}\sqrt{j!}.\label{eq:etatail}
\end{equation}
Equations \eqref{eq:E}, \eqref{eq:c} and \eqref{eq:etatail} give summable
tails for both correction series before taking the large-$N$ limit.

\newpage
\section{Closing the actual finite-time representation}\label{result}
It remains to replace $\DD e^{ish}$ by the highest-chaos marked series.
Fixed-degree bounds on $\DD^2Q_j(h)$ follow from the double commutator;
those on $\DD^2Z_j$ follow from \eqref{eq:fixed6}. Since $E_j=O(1/p)$,
spectral Cauchy--Schwarz proves
\begin{equation}
 \|\DD E_j\|_2^2\le\|E_j\|_2\|\DD^2E_j\|_2\longrightarrow0.
\end{equation}
For an all-degree bound, the nonadjacent coefficients of $Q_j$ obey
$|[x^{j-2\ell}]Q_j|\le C^\ell A^\ell\binom{j-\ell}{\ell}$.
Let $R_*=1+\sqrt{1+A}$ and $d_+=\lambda_++\lambda_+^2$.
Using the first commutator with $C$ in place of $M$ gives
\begin{equation}
 \|\DD Q_j(h)\|_2\le500\sqrt{d_+}\,j(j+1)R_*^j
                              C^{j/2}\sqrt{j!}.\label{eq:DQ}
\end{equation}
Indeed $\sum_\ell\binom{j-\ell}{\ell}A^\ell2^{j-2\ell}\le R_*^j$.
The Jacobi series therefore converges in the $\DD$ graph norm,
and fixed-degree convergence plus its summable tail yields
\begin{equation}
 \DD e^{ish}-\sum_{j\ge0}c_j(s)U_j\longrightarrow0.\label{eq:Dseries}
\end{equation}

\begin{proposition}
Define, using only the finite-parameter objects above,
\begin{align}
 \mathcal V_m(s)&=\sum_{j\ge0}c_j(s)
 \left[4\lambda_N\theta_jU_j+(\theta_j-\thN)\ehat_j\right],\label{eq:V}\\
 \widetilde B_m(s)&=\frac{m\epsilon_m}{p}C_m(1-\PP)
  e^{-ihs/2}\mathcal V_m(\alpha_ms)e^{-ihs/2}.\label{eq:Btilde}
\end{align}
Set $\widetilde B_m=0$ outside $N/4\le m\le3N/4$.
For every fixed $S$, the series is absolutely and uniformly summable,
and
\begin{equation}
 \sup_{|s|\le S}\sum_m b_m
          \|B_m(s)-\widetilde B_m(s)\|_{2,m}^2\longrightarrow0.\label{eq:full}
\end{equation}
\end{proposition}
To prove this, combine \eqref{eq:exact}, \eqref{eq:slope},
\eqref{eq:etalim} and \eqref{eq:Dseries}; $\thN U_j=\theta_jU_j$.
All factors outside the series are bounded: $m\epsilon_m/p\le2$,
$C_m\le M$, real-time phases are unitary, and pinching is contractive.
On the central annulus, both fields have uniform bounds by
\eqref{eq:E}, \eqref{eq:etatail} and the previous neutral estimate.
The binomial probability of $|m-N/2|>Xp$ is at most
$2e^{-2\lambda_-X^2}$; let $X\to\infty$ after $N$.
Outside the broad central interval, $\|B_m\|_2\le2pM$ and the
probability is at most $2e^{-N/8}$. These bounds establish \eqref{eq:full}.

\newpage
\section{Low-degree calibration and the joint-kernel task}
The exact finite formulas retain conserved terms until the final
energy pinching. Define
\begin{equation}
 b_1=(\Phi-C\theta_1)h,\quad b_2=(\Phi-C\theta_2)Z_2,
 \quad d_1=\EE\tr_m(hb_1).
\end{equation}
Equivariance makes the contraction of $h$ with $b_1$ equal to $d_1I$.
The recurrence implies
\begin{align}
 \eta_3&=-pb_1,\\
 \eta_4&=-p\{h,b_1\}-pb_2+pd_1I.\label{eq:eta4}
\end{align}
The anticommutator itself has scalar contraction $2d_1I$;
the net chaos-zero part of $\eta_4$ is $-pd_1I$.
In the central double-scale limit the earlier conventions give
\begin{equation}
 \eta_3=-F/2-4\lambda^2Cq\,h+o_{L^2}(1),\qquad
 q=e^{-4\lambda},\quad C=J_s^2/\lambda.
\end{equation}
The term proportional to $h$ is removed by final pinching, but
cannot be dropped inside a subsequent chaos operation. After pinching,
the known coefficient calibration is
\begin{align}
 (1-\PP)\eta_3&=-F/2+o_{L^2}(1),\\
 (1-\PP)\eta_4&=-2\lambda Cq^2(Y_-+Y_+)
 -(1+q)J_hF+o_{L^2}(1),
 \qquad J_hX=\{h,X\}/2.
\end{align}
These relations use the already defined marked fields of \cite{Two}.
They are coefficient checks; they do not assume interchange of
derivatives with an unevaluated limiting time kernel.

\paragraph{Implication for the research program.}
The previous $\mathsf F=A+B+o_{L^2}(1)$ and uniform second moments
now imply $\mathsf F=A+\widetilde B+o_{L^2}(1)$ on compact memory
times. Cauchy--Schwarz transfers this replacement to the two-insertion
covariance, uniformly in every real external time $T$, since
$X\mapsto e^{iTh}Xe^{-iTh}$ is unitary. All four terms
$AA,A\widetilde B,\widetilde BA,\widetilde B\widetilde B$ remain.

This is a reduction to marked inputs, not yet their spectral solution.
The outstanding calculation is the joint law of $h$, the $U_j$, their
ordered Wick descendants in \eqref{eq:etahat}, and the charged defect
with its reverse source retained. We need joint-moment formulas at every
fixed degree, uniform summation, and control of the zero-frequency
spectral mass removed by pinching. Moment convergence alone does not
justify passing an exact zero-gap projection through a limit.
The present proposition retains the finite-$N$ pinching throughout.

\newpage
\section{Checks, literature, and continuation}
\paragraph{An explicit truncation bound.}
Set $X_S=2K_*\sqrt{AM}S$ and
$C_V=4\lambda_+\sqrt{2\cdot10^6}(\sqrt{\lambda_+}+\lambda_+)+2\widehat L$.
The norm of the $j$th summand in \eqref{eq:V} is at most
$C_Ve^{2\sqrt{AM}S}(j+1)^2X_S^j/\sqrt{j!}$.
For a cutoff $J$ satisfying
\begin{equation}
 r_J=\frac{X_S}{\sqrt{J+2}}\left(\frac{J+3}{J+2}\right)^2<1,
\end{equation}
the tail after $J$ is at most
\begin{equation}
 C_Ve^{2\sqrt{AM}S}
 \frac{(J+2)^2X_S^{J+1}}{\sqrt{(J+1)!}(1-r_J)}.
\end{equation}
This intentionally loose bound certifies a truncation; it is not
an estimate of the unknown joint covariance.

\paragraph{Executed checks.}
Full-coordinate Hermite tests on $(N,p,m)=(3,1,2),(4,1,2),(4,2,3)$,
through degrees $6,5,4$, retain all $9,16,36$ Gaussian coordinates.
The 87 identity comparisons include 18 polynomial decompositions and
12 sources; maximum absolute error is $5.68\cdot10^{-15}$.
The nonzero scalar in \eqref{eq:eta4} is explicitly recovered.
The rank driver checks 12 full sixth moments at $k=3,4,6$, 38 complete
conditional distributions, 350 exact Casimir identities, and 6 independent
pure-rank mixtures. Larger-parameter rows give 54 chain and 72 conditional
\emph{upper bounds}, not direct deficit moments. All 450 slope-domination
checks pass exactly; the largest retained finite relative deviation from
the limiting slope is about $0.929$. Precision refinement (80/120 digits)
changes results by at most $1.95\cdot10^{-78}$. These checks do not replace
the uniform proof. Parameters, dependency versions and hashes are saved.

\paragraph{Existing formulas that can be reused.}
Effros--Popa \cite{EP}, Theorems 3.2--3.3, give multiblock $q$-Wick
pairing and product rules. The expected product excludes pairings
internal to a block and weights crossings; the product formula also
accounts for singletons covered by arcs. This supplies standard
combinatorics once the physical rank marks have been identified.
It does not identify $\DD$ with a spectral Wick number operator.

Berkooz et al.\ \cite{BINT}, Sections 4--5 and Appendix C.4,
provide uncrossed multiprobe kernels and the crossed four-point
$R$-matrix. Their independent probe weights make source derivatives
a possible representation of a marked insertion. Establishing that
the derivative equals our common-coupling rank mark is a separate
step; no such identification is used in the proposition.

Xu \cite{Xu}, Section 3.4, relates a chord-number generating function
to a crossed four-point function. That chord number is distinct from
$\DD=p\NN-\RR$. Direct differentiation of Eq.\ (101) yields its
$k$th moment as $(-1)^k\partial_\mu^kI|_0$; the extra $1/k!$ printed
in Eq.\ (102) does not follow from that generating-function definition.
We use the definition if pursuing this route.

\paragraph{Next calculation.}
Construct the fixed-degree joint marked contractions first, test
their lowest coefficients against the nonzero $A_0$--$B_2$ cross
measure, and only then sum the controlled series. The same-original-$J$
correlations and precise output-rank marks are the central data.
The de Sitter bulk mass/position dictionary and finite-$G$ dynamics
remain later tasks. The present result is a microscopic model
calculation, with no claim of a complete gravitational dual or of
priority over existing formulas. All artifacts remain local.

\newpage
\begin{thebibliography}{9}\small
\bibitem{Two} A. Day and Chengheng, \emph{Four-field regression and a
two-component memory field}, local stage, 7 Nov. 2026,
\texttt{susskind-memory-two-component}.
\bibitem{Uniform} A. Day and Chengheng, \emph{Uniform finite-time memory
bounds}, local stage, 7 Nov. 2026, \texttt{susskind-memory-uniform-bound}.
\bibitem{Rank} A. Day and Chengheng, \emph{Wick rank and the ordinary
propagator}, local stage, 7 Nov. 2026, \texttt{susskind-memory-wick-rank}.
\bibitem{Central} A. Day and Chengheng, \emph{Central-window sixth
moments}, local stage, 6 Nov. 2026, \texttt{susskind-memory-central-window}.
\bibitem{EP} E. G. Effros and M. Popa, \emph{Feynman Diagrams and Wick
Products Associated with $q$-Fock Space},
\href{https://arxiv.org/abs/math/0303045}{arXiv:math/0303045}.
\bibitem{BINT} M. Berkooz, M. Isachenkov, V. Narovlansky and G. Torrents,
\emph{Towards a full solution of the large N double-scaled SYK model},
\href{https://arxiv.org/abs/1811.02584}{arXiv:1811.02584}.
\bibitem{Xu} J. Xu, \emph{On Chord Dynamics and Complexity Growth in
Double-Scaled SYK},
\href{https://arxiv.org/abs/2411.04251}{arXiv:2411.04251v3}.
\end{thebibliography}
\end{document}
